\section{The ordinary rank-two factor of a Case III seed}

% BEGIN R66P1
\paragraph{R66P1. A rational prime splitting the entire frame and its auxiliary character.}
Fix the source Case III datum $D$ over $O=c.O$ and the initial
tower $T_0$ as in $\mathrm{RhoExP}$. Write $G=G_{\mathbb Q}$,
$F_0=T_0.F$, $E=\overline{\mathbb Q}_\ell$ and
$\varphi=\xi\psi^j$, where $j=1$ for the false label and
$j=-1$ for the true label. We use the indicated coefficient
embedding $O\subset E$ on all character values.

Both factors of $\varphi$ factor through a cyclotomic field of level
prime to $\ell$. Indeed, the Case III host presents $\xi$ as the
product of Dirichlet characters of levels $p_1$ and
$2^{n_{\rm III}}$, and $\psi$ as a Dirichlet character of level
$N_H=\operatorname{lcm}(2^{\mathrm{ex}+2},q)q_1d$.
The excluded primes in the host include every odd prime divisor of
$qp_1$. Also $\ell-1\mid2e$ with $0<e<t$, whereas
$q_1,d>2t+1$ are primes. Thus $\ell$ divides none of these
levels. Take $N=p_1 2^{n_{\rm III}}N_H$. An element fixing
$\mu_N$ pointwise fixes the roots of unity defining both factors,
so $\varphi$ kills $G_{\mathbb Q(\mu_N)}$ and has finite image.

The subgroup
\[
       U=G_{F_0}\cap\ker\bar r_D\cap\ker\omega
\]
is open and has abelian quotient. In fact $F_0/\mathbb Q$ is
finite abelian, the diagonal residual representation has continuous
finite image, and $\omega$ is the finite continuous lift of the
residual cyclotomic character. By Kronecker--Weber the fixed field
of $U$ lies in $\mathbb Q(\mu_M)$ for some $M$. Choose a prime
$w>N\ell$ congruent to $1$ modulo $MN\ell$ using Dirichlet's
theorem. Then all decomposition groups at $w$ lie in $U$, since
$w$ splits completely in $\mathbb Q(\mu_M)$. This is precisely
$T_0.\mathrm{Splits}(w)$. Also every element acting as
$\zeta\mapsto\zeta^w$ on $\mu_N$ fixes $\mu_N$, so it is
killed by $\varphi$. This verifies both clauses of $\mathrm{SplitAt}$.

Send a generator of $(\mathbb Z/w\mathbb Z)^\times$ to a primitive
$\ell$th root of unity in $E$. This defines a character because
$\ell\mid w-1$. Composing with the surjective action of $G$ on
$\mu_w$ gives a nontrivial $\theta$ of exact order $\ell$ killing
$G_{\mathbb Q(\mu_w)}$. We have $N>0$, $\ell\nmid N$, $w\nmid N$,
and $\ell$ is odd. These are all the hypotheses collected in
$\mathrm{SlopeTestRemainingChoices}(N,w,\theta)$.
% END R66P1

% BEGIN R66P2
\paragraph{R66P2. Parity and the two residual distinctions.}
Complex conjugation $c_\infty$ has order two and $\psi$ has odd
order in Case III, hence $\psi(c_\infty)=1$. As $\xi(c_\infty)=-1$,
we have $\varphi(c_\infty)=-1$. The host parity
$\psi(-1)=(-1)^e$ of the underlying Dirichlet character gives
$e$ even by the same argument.

Write $\varepsilon$ for the full cyclotomic character, put
$\Phi=\varepsilon^{k-1}\varphi$ and $\Psi=\varepsilon^e\psi$,
and suppose $k\ge1$ and $k-1\equiv e\pmod{\ell-1}$.
Bars denote residual characters. Since $\bar\varepsilon^{2e}=1$,
we have
\[
   \bar\Phi=\bar\xi\,\bar\Psi^j
       =\bar\chi_0/\bar\chi_3 .
\]
For the last equality, put $\alpha=\varepsilon\omega^{-1}\xi$,
so $\bar\alpha=\bar\xi$. In the false labelling the ratio
is $\bar\alpha\bar\Psi$, and in the true labelling it is
$\bar\alpha/\bar\Psi$.

Let $g_0$ be the rational image of the initial tower mark, and write
$b_i=\bar\chi_i(g_0)$. The four $b_i$ are pairwise distinct.
Moreover $b_0/b_1$ is the value of the character of the supplied
quadratic field and is $\pm1$. Thus $b_0=-b_1$. If
$(b_0/b_3)^2=1$, then $b_0=b_3$ or $b_1=b_3$, both impossible.
This proves $\bar\Phi(g_0)^2\ne1$. At the fixed rational place
of $D$ above $\ell$, the specified Frobenius lift $D.\mathrm{sg}(\ell)$
in the decomposition group satisfies
$(\bar\xi\,\bar\Psi^j)(D.\mathrm{sg}(\ell))\ne1$ by $D.\mathrm{Good}$.
Conjugation of places preserves this character value. Thus at every
place $A$ above $\ell$ the rational decomposition group $D_A$ has
some $s_A$ with $\bar\Phi(s_A)\ne1$.
% END R66P2

% BEGIN R66P3
\paragraph{R66P3. A sufficiently large ordinary weight and irreducibility on open subgroups.}
Apply the stated hypothesis $\mathrm{PropEGP}$ with the data
of R66P1. Its hypothesis $\varphi\ne1$ follows from R66P2.
It yields a threshold $k_0$. Given the requested lower bound $B$, take
\[
           k=e+1+(\ell-1)(B+k_0+2).
\]
Then $k\ge B$, $k\ge k_0$, $k>1$, and $k-1\equiv e\pmod{\ell-1}$.
This weight is odd, so its required sign is exactly
$\varphi(c_\infty)=(-1)^k$. The hypothesis gives the same
$\rho:G\to\mathrm{GL}_2(E)$ with $\mathrm{EisRep}(\iota,\varphi,\theta,k,\rho)$
and $\mathrm{PR}(\iota,\rho)$. The latter includes a regular
algebraic cuspidal attachment satisfying the source's $\mathrm{PlainR}$.

We strengthen the irreducibility given by $\mathrm{EisRep}$.
Suppose $\rho$ were reducible on an open subgroup. Its open normal
core $H$ would then preserve a line $V$. The distinguished inertia
element $s_0$ over $\ell$ has full cyclotomic value of infinite
order. A positive power $s_0^n$ belongs to $H$. The ordinary
inertial characters of $\rho$ there are $\varepsilon^{k-1}$ and
$1$, so this element has two distinct eigenvalues. Every $G$-translate
of $V$ is $H$-invariant by normality and must be one of these two
eigenlines. Both lines occur, since $\rho$ is irreducible over $G$.
Thus $G$ acts transitively on a pair of lines. In their basis each
element interchanging them has trace zero. Since the characteristic
polynomial of each $\rho(g)$ reduces to
$(X-1)(X-\bar\Phi(g))$, the trace formula gives
$\bar\Phi(g)=-1$ for every interchanging element. Multiplying any
element in the stabilizer by a fixed interchanging element then
gives $\bar\Phi(g)=1$ on the stabilizer. This forces
$\bar\Phi^2=1$ everywhere, contradicting the value at $g_0$ in
R66P2. This proves $\mathrm{IrrOpen}(\rho)$.
% END R66P3

% BEGIN R66P4
\paragraph{R66P4. An integral frame diagonal at the prescribed rational mark.}
In the field $E$, call two scalars congruent if their difference
has norm $<1$. The residual characteristic polynomial of $\rho(g_0)$
has distinct roots $\bar\Phi(g_0),1$. Its characteristic-zero
roots are integral (the polynomial is monic with integral
coefficients), and their reductions are these two roots. Diagonalize
in this order, with eigenvalues $u,v$ and $|u-v|=1$. Write
$\rho(g)=\left(\begin{smallmatrix}a_g&b_g\\c_g&d_g\end{smallmatrix}\right)$
in that basis. The traces at $g$ and $g_0g$ give
\[
    a_g+d_g\equiv\Phi(g)+1,\qquad
    u a_g+v d_g\equiv\Phi(g_0)\Phi(g)+1 .
\]
Subtract the same two expressions evaluated at $(\Phi(g),1)$.
The coefficient matrix has integral inverse, so
$a_g\equiv\Phi(g)$ and $d_g\equiv1$ even before assuming
integrality of the off-diagonal entries. Taking entry $(0,0)$ in
the product for $gh$ gives $b_gc_h=a_{gh}-a_ga_h\equiv0$ for
every $g,h$. Some $c_h$ is nonzero, since otherwise the first
basis line would be invariant. Continuity and compactness give a
maximum of $|c_h|$ attained at $h_*$. Diagonal conjugation
scales every $c_h$ by $c_{h_*}^{-1}$ and every $b_g$ by $c_{h_*}$.
We now have
\[
    a_g\equiv\Phi(g),\quad d_g\equiv1,\quad
    |b_g|<1,\quad |c_g|\le1,\quad c_{h_*}=1,
    \qquad b_{g_0}=c_{g_0}=0 .
\]
This is a single continuous representation isomorphic to $\rho$
in the required order, with a unit lower-left entry somewhere on $G$.
% END R66P4

% BEGIN R66P5
\paragraph{R66P5. One finite coefficient ring for every ordinary flag.}
The matrices in the frame of R66P4 are defined over a single finite
extension $K/\mathbb Q_\ell$. To see this, the finite subextensions
inside $E$ have a countable cofinal family. Indeed
$\overline{\mathbb Q}$ has a countable dense image in $E$ under
a fixed embedding (approximate a separable minimal polynomial by
rational polynomials and use continuity of its simple roots).
Krasner's lemma then places each algebraic $z$ in a field generated
by one sufficiently close element of that image. Taking finite
composita gives the claimed family. For each such $K_i$, the
elements whose matrices and inverse matrices have all entries in
$K_i$ form a closed subgroup of $G$ since the finite field
extension $K_i$ is complete. These subgroups cover $G$. By the
Baire theorem for compact Hausdorff groups, one has interior and is
therefore open. Adjoining its finitely many coset representatives'
matrix entries gives the desired $K$.

At each place $A$ over $\ell$ the ordinary input gives a line
for $D_A$ whose character $\kappa_0$ reduces to $\bar\Phi$.
The quotient character is $\kappa_1$. Write a primitive vector
on this line, by scaling to maximum coordinate norm one. If its
first coordinate had norm $<1$, the line would reduce to the
second basis axis. The integral matrices of R66P4 act on that
axis in the reduction by $1$, contradicting
$\bar\kappa_0(s_A)=\bar\Phi(s_A)\ne1$. The first coordinate
is therefore a unit and the line has generator $(1,x_A)^t$ with
$|x_A|\le1$. For $g\in D_A$ one obtains
\[
 a_g+b_gx_A=\kappa_0(g),\qquad
 c_g+d_gx_A=(a_g+b_gx_A)x_A,\qquad
 d_g-b_gx_A=\kappa_1(g).
\]
The last equality also follows by taking the trace. On inertia
the first and last characters are exactly $\varepsilon^{k-1}$
and $1$.

These slopes all belong to $K$. At $s_A$ we have $|a-d|=1$ and
$|b|<1$, with $a,b,c,d$ denoting that single integral matrix.
The displayed flag identity makes $x_A$ an integral fixed point of
$y\mapsto(c-by^2)/(a-d)$. This maps the closed unit ball of $K$
to itself and contracts distances by at most $|b|$. Iteration
from zero converges in $K$. The same inequality on the unit ball
of $E$ gives uniqueness, hence its limit is $x_A$. Set $x_A=0$
at places away from $\ell$.

Take the compositum of $K$ with the coefficient fields of $c$ and
the prescribed $c_0$. Its ring of integers $O_1$ gives a finite
coefficient datum containing both $O$ and $c_0.O$. Every matrix
and every $x_A$ takes values in $O_1$. Its maximal ideal consists
of the elements of norm $<1$, and its powers are the balls of radii
$|\varpi|^n$ for a uniformizer $\varpi$. Thus all matrix entries
are adically continuous over $O_1$. Entrywise inclusion into $E$
recovers the conjugate of $\rho$ in R66P4.
% END R66P5

% BEGIN R66P6
\paragraph{R66P6. The normalized modular model and the native supplier.}
Denote the representation over $O_1$ by $\rho_1$. Its
characteristic-zero isomorphism, entrywise continuity, residual
diagonal $(\bar\Phi,1)$, upper-right congruence, lower-left unit
witness and diagonal normalization at $g_0$ are R66P4--5.
On $\ker\theta$ its determinant is the inclusion of $\Phi(g)$ since
$\mathrm{EisRep}$ has determinant $\varepsilon^{k-1}\varphi\theta$.
At every place away from $\ell$ the inertia type in that same input
is conjugate to $\operatorname{diag}(1,\varphi\theta)$. Thus
$\rho_1(g)=1$ for every such inertia element with
$\varphi(g)=\theta(g)=1$.
The flag identity and the two exact inertia identities at all places
over $\ell$ are those of R66P5. Together these verify all nine
fields of $\mathrm{Mod4}$ for the original $g_0,k,\theta$ and $\rho$.

The resulting $w,\theta,k,\rho,O_1,\rho_1,x$ satisfy
$\mathrm{RhoExP}$ for the original complex comparison, the
arbitrary bound $B$ and the independently prescribed coefficient
datum $c_0$. R66P1 supplies the splitting and
$\mathrm{KillsMu}(w,\theta)$, R66P3 supplies the weight, its
congruence, $\mathrm{IrrOpen}$ and the cuspidal attachment with
$\mathrm{PR}$, and R66P5 supplies both coefficient inclusions.
This proves the implication with the Eisenstein-family hypothesis
$\mathrm{PropEGP}$ appearing in its definition.
% END R66P6
